\documentclass[11pt]{article}
\usepackage[a4paper,margin=22mm]{geometry}
\usepackage[no-math]{fontspec}
\setmainfont{DejaVu Serif}
\usepackage{amsmath,amssymb,amsthm,mathtools,graphicx,xcolor,hyperref,xurl,xspace,ifthen}
\setlength{\emergencystretch}{4em}
\SetMathAlphabet{\mathbf}{normal}{OT1}{cmr}{bx}{n}
\SetMathAlphabet{\mathbf}{bold}{OT1}{cmr}{bx}{n}
\newtheorem{theo}{Theorem}[section]
\newtheorem{lemm}[theo]{Lemma}
\newtheorem{theorem}[theo]{Theorem}
\newtheorem*{remas*}{Remarks}
\newtheorem{remas}[theo]{Remarks}
\newtheorem*{exam*}{Example}
\newtheorem{ques}[theo]{Question}
\providecommand{\og}{«\nobreakspace}
\providecommand{\fg}{\nobreakspace»}
\newtheorem{lem}[theo]{Lemma}
\newtheorem{lemma}[theo]{Lemma}
\newtheorem{prop}[theo]{Proposition}
\newtheorem{coro}[theo]{Corollary}
\newtheorem{corol}[theo]{Corollary}
\newtheorem{conj}[theo]{Conjecture}
\newtheorem{Proposition}[theo]{Proposition}
\newtheorem{theoreme}[theo]{Theorem}
\theoremstyle{definition}
\newtheorem{defi}[theo]{Definition}
\newtheorem{exam}[theo]{Example}
\newtheorem{exem}[theo]{Example}
\newtheorem{remark}[theo]{Remark}
\newtheorem{example}[theo]{Example}
\newtheorem{rema}[theo]{Remark}
\newtheorem{nota}[theo]{Notation}
\newtheorem*{rema*}{Remark}
\newtheorem*{defi*}{Definition}
\newtheorem*{theo*}{Theorem}
\newtheorem*{lemm*}{Lemma}
\newtheorem*{prop*}{Proposition}
\newcommand{\enoncename}{}
\newtheorem{corpusenonce}[theo]{\enoncename}
\newenvironment{enonce}[2][]{\renewcommand{\enoncename}{#2}\begin{corpusenonce}}{\end{corpusenonce}}
\newenvironment{enonce*}[2][]{\par\medskip\noindent\textbf{#2.}\itshape}{\par\medskip}
\newenvironment{proofc}[1][\proofname]{\begin{proof}[#1]}{\end{proof}}
\newcommand{\Mc}{\mathcal{M}}
\newcommand{\pointir}{.\ ---\ }
\providecommand{\mathof}[1]{\mathrm{#1}}

\numberwithin{equation}{section}
\numberwithin{figure}{section}
\numberwithin{table}{section}
\def\endappendix{}
%~Mouliné par MaN_auto v.0.21.4 2019-04-09 15:38:57


\newfontfamily\CorpusDoubleStroke{Noto Sans Math}
\newcommand{\mathds}[1]{{\ifthenelse{\equal{#1}{1}}{\mathord{\text{\CorpusDoubleStroke\char"1D7D9}}}{\mathbb{#1}}}}









%\DeclareGraphicsRule{*}{mps}{*}{}

\DeclareMathOperator{\supp}{supp}
\DeclareMathOperator{\Reop}{Re}
\renewcommand{\Re}{\Reop}

\newcommand*\bbbone{\mathds{1}}

\newcommand{\CI}{{\mathcal C}^\infty }
\newcommand{\CIc}{{\mathcal C}^\infty_{\mathrm{c}} }
\newcommand{\CC}{{\mathbb C}}
\newcommand{\NN}{{\mathbb N}}
\newcommand{\R}{{\mathbb R}}
\newcommand{\TT}{{\mathbb T}}
\newcommand{\ZZ}{{\mathbb Z}}



\newcommand{\defeq}{\stackrel{\mathrm{def}}{=}}

\newcommand{\loc}{\mathrm{loc}}

\newcommand*\Z{\mathbb Z}
\newcommand*\N{\mathbb N}

%\newcommand*\squarebox[1]{\hbox to #1{\hfill\vbox to #1{\vfill}}}

\newcommand*{\dd}{\mathrm{d}}
\newcommand*{\dx}{\dd x}
\newcommand*{\dy}{\dd y}
\newcommand*{\dz}{\dd z}
\newcommand*{\ds}{\dd s}
\newcommand*{\dt}{\dd t}
\newcommand*{\dxi}{\dd \xi}
\newcommand*{\dmu}{\dd \mu}
\newcommand*{\dwmu}{\dd \widetilde\mu}
\newcommand*{\dwnu}{\dd \widetilde\nu}

\graphicspath{{./figures/}}

\newcommand*{\mk}{\mkern -1mu}
\newcommand*{\Mk}{\mkern -2mu}
\newcommand*{\mK}{\mkern 1mu}
\newcommand*{\MK}{\mkern 2mu}

\hypersetup{urlcolor=purple, linkcolor=blue, citecolor=red}

\newcommand*{\romanenumi}{\renewcommand*{\theenumi}{\roman{enumi}}}
\newcommand*{\Romanenumi}{\renewcommand*{\theenumi}{\Roman{enumi}}}
\newcommand*{\alphenumi}{\renewcommand*{\theenumi}{\alph{enumi}}}
\newcommand*{\Alphenumi}{\renewcommand*{\theenumi}{\Alph{enumi}}}

\newcommand{\llbracket}{\mathopen{[\mkern -2.7mu[}}
\newcommand{\rrbracket}{\mathclose{]\mkern -2.7mu]}}
\let\oldexists\exists
\renewcommand*{\exists}{\mathrel{\oldexists}}
\let\oldforall\forall
\renewcommand*{\forall}{\mathrel{\oldforall}}

\newenvironment{remasc*}{\let\pointir\relax \begin{remas*}}{\end{remas*}}
%\renewcommand\thefootnote{\ddag}
\newcommand*{\tq}{\,|\,}
\newcommand*{\Tq}{\,\middle|\,}

\begin{document}
\section*{\raggedright 2688: Arbitrary-time Schrödinger observability from nonzero nonnegative L-four weights on rational rectangular two-tori}
\noindent\textbf{Author:} Nicolas Burq; Maciej Zworski\par\noindent\textbf{Source:} Rough controls for Schrödinger operators on 2-tori. Annales Henri Lebesgue 2 (2019), publisher version, DOI 10.5802/ahl.19. \url{https://ahl.centre-mersenne.org/item/AHL_2019__2__331_0/}
\par\noindent\textbf{License:} CC BY 4.0, \url{https://creativecommons.org/licenses/by/4.0/}. The original publisher PDF has the explicit license notice. See the saved source/license records.
\par\noindent\textbf{Editorial material note (not author text):} The selected problem is Theorem 1.3 restricted to the nonzero rational rectangular torus parameter and a nonzero nonnegative L4 observation weight, for every positive time. The complete author body retains all local semiclassical defect-density, rational-direction, weighted one-dimensional, dyadic decomposition/commutator, low-frequency compactness and unique-continuation arguments. The explicitly published mixed-norm estimate and Jakobson limit theorem remain named prerequisites. Irrational-torus extension, general complex weights, controllability and damping assertions are unselected context. All proof text and formulas below are editable author TeX. Mathematical wording, language and proof order are retained. Only the article class, title matter and bibliography presentation are adapted for independent compilation; no proof has been generated or repaired.
\par\noindent\textbf{Editorial difficulty:} H2: The authored proof extracts L-two densities of semiclassical defects, isolates rational directions, reduces rough observation to one-dimensional dynamics, and removes low-frequency error by compactness and unique continuation. This rough-control microlocal mechanism exceeds ordinary qualifying PDE.
\medskip\hrule\medskip
\noindent\textbf{Author statement, complete proof and local supporting text follow in source order.}
\section{Introduction}\label{in}

The purpose of this paper is to investigate the general question of control theory with localized control functions. When the localization is performed by a continuous function, the question is completely settled for wave equations~\cite{BLR,BG} and well understood for Schrödinger equations on tori~\cite{Ha, Ja, Ko, BZ4, AM}.

In this paper we localize only to sets of \emph{positive measure} or more generally use control functions in $ L^2 $. The understanding is then much poorer and only partial results are available even for the simpler case of wave equations~\cite{Bu16,Bu17}. Using the work with Bourgain~\cite{BBZ} and~\cite{BZ4} we completely settle the question for Schrödinger equation on the two dimensional torus taking advantage, as in previous papers, of the particular simplicity of the dynamical structure.


To state the control result consider
\begin{equation}\label{eq:con}
\begin{gathered} \TT^2 : = \R^2 / \ZZ \times \gamma \ZZ \,, \ \ \gamma \in \R \setminus \{
0 \}\,, \ \ a \in L^2 (\TT^2), \\
(i \partial_t + \Delta) u (t, z) = a(z) 1_{(0,T)}f \,, \ \ \ u (0, z) =u_0 (z),
\end{gathered}
\end{equation}
where $a $ is a localisation function and $f$ a control. From~\cite[Proposition~2.2]{BBZ} (see Theorem~\ref{th.BBZ} below) we know that for $f \in L^{4}(\TT^2 ; L^2(0,T))$ (so that $af \in L^{4/3}(\TT^2 ; L^2(0,T))$), and any $u_0 \in L^2(\TT^2)$, there exists a unique solution
\[
u \in C^0 ([0,T]; L^2(\TT^2)) \cap L^4 (\TT^2; L^2(0,T)).
\]

A classical question of control is to fix $ a $ and ask for which $u_0 \in L^2$ does there exist a control $f$ such that the solution of~\eqref{eq:con} satisfies $u |_{t>T} =0$? We show that on $\mathbb{T}^2$ it is always the case as soon as $a \in L^2$ is non-trivial:

\begin{theo}\label{th.2}
Let $a \in L^2(\TT^2), \| a \|_{L^2} >0$ and $ T > 0 $. Then for any $u_0\in L^2(\TT^2)$ there exists $f \in L^4(\TT^2 ; L^2(0,T))$ such that the solution $u$ of~\eqref{eq:con} satisfies $u |_{t=T} =0$.
\end{theo}

The next result shows that adding an $ L^2 $ damping term results in exponential decay:

\begin{theo}\label{th.3}
For $a\in L^2(\TT^2)$, $ a \geq 0 $, $ \| a \|_{L^2} > 0$, there exist $C, c>0$ such that for any $u_0 \in L^2(\TT^2)$, the equation
\begin{equation}\label{damped}
(i \partial_t + \Delta + ia) u=0, \qquad u |_{t=0} = u_0,
\end{equation}
has a unique global solution $ u \in L^\infty(\R; L^2(\TT^2))\cap L^4(\TT^2; L^2_{\loc}(\R)) $ and
\begin{equation}\label{dampedbis}
\| u\|_{L^2(\TT^2)} (t) \leq C e^{-ct} \| u_0 \|_{L^2(\TT^2)}.
\end{equation}
\end{theo}

As shown in Section~\ref{HUM} both results follow from an the observability estimate. We should think of $ a $ in Theorem~\ref{th.2} as $ W^2 $ where $ W $ appears in the following statement:

\begin{theo}\label{th.1}
Suppose that $ W \in L^4 (\TT^2) $, $ \| W \|_{L^4} >0 $. Then for any $ T > 0$ there exists $ K $ such that for $ u \in L^2 (\TT^2) $,
\begin{equation}\label{eq:th1}
\| u \|_{ L^2 (\TT^2) } \leq K \| W e^{ it \Delta} u \|_{ L^2 ((0,T)_t \times \TT^2) }.
\end{equation}
\end{theo}

To keep the paper easily accessible we present proofs in the case when $ \gamma \in \mathbb Q $ in~\eqref{eq:con}. Irrational tori require a more complicated reduction to rectangular coordinates (see~\cite[Lemma~2.7 and Figure~1]{BZ4} but the modification can be done as in that paper). The crucial~\cite[Proposition~2.2]{BBZ} is valid for all tori. Another approach to treating (higher dimensional) irrational tori can be found in the work of Anantharaman--Fermanian-Kammerer--Macià, see~\cite[Corollary~1.19, Theorem~1.20]{AFM}.

Since, as is already clear, \cite[Proposition~2.2]{BBZ} plays a central role in many proofs we recall it in a version used here:

\begin{theo}\label{th.BBZ}
Let $ T > 0 $. There exists $ C = C_T $ such that for
\[
u_0\in L^2(\TT^2), \ \ \ f \in L^{\frac 4 3 } (\TT^2 ; L^2(0, T)),
\]
the solution to $ (i \partial_t + \Delta) u =f$, $ u|_{t=0} = u_0$, satisfies
\[
\| u\|_{L^\infty ((0, T) ; L^2 (\TT^2) \cap L^4(\TT^2; L^2((0,T)))}\\
\leq C \left(\|u_0\|_{L^2(\TT^2)} + \|f\|_{ L^1 ((0,T) ; L^2 (\TT^2)) + L^{\frac 4 3 } (\TT^2 ; L^2(0, T))}\right).
\]
\end{theo}

\begin{remasc*}\ 
\begin{enumerate}
\item \label{rem1} Theorem~\ref{th.1} is equivalent to the same statement with $ W \in L^\infty (\TT^2)$ (by replacing $ W \in L^4 $ by $ \bbbone_{ |W | \leq N } W \in L^\infty $ with $ N$ sufficiently large noting\linebreak that $ \| W v \|_{ L^2 (\TT^2) } \geq \| W \bbbone_{ |W | \leq N } W v \|_{ L^2 (\TT^2) } $ and that $ \|W \|_{L^4} > 0 $ implies \linebreak$\| W \bbbone_{ |W | \leq N } \|_{L^\infty}> 0 $ if $ N $ is large enough. Our formulation is more consistent with the statements of Theorems~\ref{th.2} and~\ref{th.3}.

\item \label{rem2} For rational tori and for $ T > \pi $, Theorem~\ref{th.1}, and by Proposition~\ref{p:hum} below, Theorems~\ref{th.2} and~\ref{th.3}, follow from the results of Jakobson~\cite{Jak}. That is done by using the complete description of microlocal defect measures for eigenfunctions of $
\mathcal \R^2 / 2 \pi \ZZ^2 $. We explain this in detail in the appendix.

\item \label{rem3} The starting point of~\cite{Jak} and~\cite{BBZ} was the classical inequality of Zygmund:
\begin{multline}\label{eq:zyg0}
\forall \lambda \in \N,\ \left\| \sum_{{ |n|^2= \lambda}} c_ne^{in \cdot z} \right\|_{L^4(\TT^2_z)}^2 \leq \frac{\sqrt 5}{ 2 \pi}
\sum_{{ |n|^2= \lambda}} |c_n|^2, \\ z \in \TT^2 = \R^2 / 2 \pi \ZZ^2, \ \ n \in \ZZ^2.
\end{multline}
In particular for $ \TT^2 = \R^2/ 2 \pi \ZZ^2 $, we easily see how the homogeneous part ($f=0$) in Theorem~\ref{th.BBZ} follows from~\eqref{eq:zyg0}. For that put $ u = \sum_{\lambda } u_\lambda$, $ u_\lambda = \sum_{ _{ { |n|^2= \lambda}} } c_n e^{ i n \cdot x} $. Then, using~\eqref{eq:zyg0} in the third line,
\begin{align*}
\| e^{ it \Delta } u \|_{ L^4 (\TT^2, L^2 ((0, 2 \pi))) }^4 & =
\int_{ \TT^2 } \left(\int_{0}^{2\pi } \left| \sum_{ \lambda } e^{ it \lambda} u_\lambda (z) \right|^2 \dt \right)^2 \dz \\
& = (2 \pi)^2 \int_{ \TT^2 } \left(\sum_\lambda | u_\lambda (z) |^2 \right)^2 \dz \\
& = (2 \pi)^2\int_{\TT^2 } \sum_{ \lambda, \mu} |u_\lambda (z) |^2 |u_\mu (z) |^2 \dz \\
&\leq (2 \pi)^2\sum_{ \lambda, \mu } \| u_\lambda \|_{L^4}^2 \| u_\mu \|_{L^4}^2\\
& \leq 5 \sum_{ \lambda, \mu} \| u_\lambda \|_{L^2}^2 \| u_\mu \|_{L^2}^2 = 5 \left(\sum_\lambda \| u_\lambda \|_{L^2}^2 \right)^2 = 5 \| u \|_{L^2}^4.
\end{align*}
Generalizations for the time dependent Schrödinger equation in higher dimensions were obtained by Aïssiou--Jakobson--Macià~\cite{AJM}.

\item \label{rem4} Other than tori, the only other manifolds for which~\eqref{eq:th1} is known for \emph{any}\linebreak non-trivial continuous $ W $ are compact hyperbolic surfaces. That was proved by Jin~\cite{Ji} using results of Bourgain--Dyatlov~\cite{BD} and Dyatlov--Jin~\cite{DJ}.
\end{enumerate}
\end{remasc*}

\paragraph*{Acknowledgements}
This research was partially supported by Agence Nationale de la Recherche through project ANAÉ ANR-13-BS01-0010-03 (NB), by the NSF grant DMS-1500852 and by a Simons Fellowship (MZ). MZ gratefully acknowledges the hospitality of Université Paris-Sud during the writing of this paper. We would also like to thank Alexis Drouot, Semyon Dyatlov and Fabricio Macià for helpful comments on the first version. We are also grateful to the referee for the careful reading of the first version of the paper and for the many helpful comments.


\section{Semiclassical observability}

We follow the strategy of~\cite{BZ4} and~\cite{BBZ} and first prove a semiclassical observability result. For that we define
\begin{equation}\label{eq:Pih}
\Pi_{h, \rho} (u_0) := \chi \left(\frac{ - h^2\Delta - 1} {\rho} \right) u_0\,, \ \ \ \rho > 0 \,,
\end{equation}
where $ \chi \in \CIc ((- 1, 1))$ is equal to $ 1 $ near $ 0 $. With this notation the main result of this section is
\begin{prop}\label{p:semi}
Suppose that $ a \in L^2 (\TT^2) $, $ a \geq 0 $, $ \| a \|_{L^2} > 0 $. For any $ T > 0 $ there exist $ K $, $ \rho_0 > 0 $ and $ h_0 > 0 $ such that for any $ u_0 \in L^2 (\TT^2) $,
\begin{equation}\label{eq:semi}
\| \Pi_{h, \rho} u_0 \|_{ L^2 }^2 \leq K \int_0^T \!\!\int_{ \TT^2} a (z) | e^{ i t \Delta} \Pi_{h, \rho} u_0 (z) |^2 \dz \dt,
\end{equation}
for $ 0 < \rho < \rho_0 $ and $ 0 < h < h_0 $.
\end{prop}

The proof of the Proposition proceeds by contradiction: if~\eqref{eq:semi} does not hold then there exists $ T > 0 $ such that for any $ n \in \NN $ there exist $ 0 < h_n < 1/n $, $ 0 < \rho_n < 1/n $ and $ u_n \in L^2 $ for which
\begin{equation}\label{eq:contr}
1 = \| u_n \|_{ L^2 }^2 > n \int_0^T \!\!\int_{ \TT^2} a (z) | e^{ i t \Delta} u_n (z) |^2 \dz \dt, \ \ u_n = \Pi_{ h_n, \rho_n } u_n.
\end{equation}
We will use semiclassical limit measures associated to subsequences of $ u_n$'s.

\subsection{Semiclassical limit measures}

Each sequence $ u_n (t) := e^{ i t \Delta } u_n $, is bounded in $L^2_{\loc} (\mathbb{R} \times \TT^2)$. After possibly choosing a subsequence, $ u_n $'s define a semiclassical defect measure $\mu $ on $\R_t \times T^* (\TT^2_z)$ such that for any function $\varphi \in {\mathcal C}^0_{\mathrm{c}} (\R_t)$ and any $A\in \CIc (T^*\TT^2_z)$, we have
\begin{equation}\label{eq:mu}
\langle \mu, \varphi(t) A(z, \zeta)\rangle = \lim_{n\rightarrow \infty} \int_{\R_t } \varphi(t)
\langle A(z, h_{n} D_z) u_{n} (t), u_n (t)
\rangle_{ L^2 (\TT^2) } \dt \,.
\end{equation}

The measure $ \mu $ enjoys the following properties:
\begin{gather}\label{eq:propmu}
\begin{gathered}
\!\mu ((t_0, t_1) \Mk\times\Mk T^*\TT^2_z) \mk=\mk t_1 \mk-\mk t_0, \ \
\supp \mu \mk\subset\mk \Sigma \mk:=\mk \{(t,z, \zeta) \mk\in\mk \R_t \Mk\times\Mk \TT^2_z \Mk\times\Mk \R^2_\zeta \tq |\zeta|\mk=\mk 1\},\! \\
\partial_s \int_{\R} \int_{ T^* \TT^2 } \varphi(t) A (z + s \zeta, \zeta) \dmu = 0,
\ \ \varphi \in {\mathcal C}^0_{\mathrm{c}} (\R_t), \ A \in \CIc (T^*\TT^2_z),
\end{gathered}
\end{gather}
see~\cite{Ma} for the derivation and references.

We have an additional property which follows from an easy part of Theorem~\ref{th.BBZ} (in the rational case related to the Zygmund inequality~\eqref{eq:zyg0}; see also Aïssiou--Jakobson--Macià~\cite[Theorem~1.2]{AJM}): for any $ \tau \geq 0 $ there exists $ m_\tau
\in L^2 (\TT^2) $ such that for all $ f \in \mathcal C (\TT^2) $
\begin{equation}\label{eq:mT}
\int_{0}^\tau \!\!\int_{ T^* \TT^2 } f (z) \dmu (t, z, \zeta) =
\int_{\TT^2 } m_\tau (z) f (z) \dz.
\end{equation}
In fact, Theorem~\ref{th.BBZ} shows that
\begin{equation}\label{eq:UnT}
U_n^\tau (z) := \int_{0}^\tau | u_n (t, z) |^2 \dt
\end{equation}
satisfies
\begin{equation}\label{eq:prop22}
\| U^\tau_n \|_{ L^2 (\TT^2_z) } = \| u_n (t, z) \|_{ L^4 (\TT^2_z, L^2 ((0, \tau)) }^2 \leq C \| u_n \|_{L^2 (\TT^2) }^2 = C.
\end{equation}
But then, after passing to a subsequence, $ U_n^\tau $ converges weakly to $ m_\tau \in L^2 (\TT^2) $. Since
\[
\int_{0}^\tau \int_{ T^* \TT^2 } f (z) \dmu (t, z, \zeta) = \lim_{ n \to \infty } \int_{0}^\tau \int_{\TT^2} f (z) U_n^\tau (z) \dz,
\]
this proves~\eqref{eq:mT}.

Using~\eqref{eq:contr} and then the fact that $ U_n^T $ given in~\eqref{eq:UnT} converge weakly to $ m_T $ in $ L^2$ we obtain
\begin{equation}\label{eq:l1}
\int_{\TT^2 } a (z) m_T (z) \dz = 0.
\end{equation}

The next lemma shows that our measure has most of its mass on the set of rational directions:

\begin{lemm}\label{l:2}
Suppose that $\mu $ is defined by $ u_n$ satisfying~\eqref{eq:contr}. For $ m \in \NN $ define,
\[
W^m \!:=\! \left\{ (z, \zeta) \in T^* \TT^2 \Tq \zeta = \frac{ (p, q) }{
\sqrt{ p^2 + q^2 }}, \, \max(|p|, |q|) \leq m, (p, q) \in \Z^2, \, \gcd(p,q) = 1 \right\}\!,
\]
its complement, $ W_m := \complement W^m $, and a measure $ \widetilde \mu_T $ on $ T^* \TT^2 $:
\begin{equation}\label{eq:defwm}
\int_{T^* \TT^2} A (z, \zeta) \dwmu_T := \int_0^T \int_{T^* \TT^2 } A (z, \zeta) \dmu (t, z, \zeta), \ \ A \in \CIc (T^* \TT^2).
\end{equation}
Then,
\begin{equation}\label{eq:l2}
\forall \epsilon > 0 \ \exists m \ \text{ such that } \
\widetilde \mu_T (W_m) < \epsilon.
\end{equation}
\end{lemm}

\begin{proof}
We choose
\begin{equation}\label{eq:aj}
a_j \in \CI (\TT^2), \ \ a_j \geq 0, \ \
\lim_{ j \to \infty } \| a - a_j\|_{L^2 } = 0.
\end{equation}
\begin{equation}
Then\label{eq:muTaj}
\int_{T^* \TT^2 } a_j (z) \dwmu_T (z,\zeta) = \int_{ \TT^2 } (a_j (z) - a (z)) m_T (z) \dz = \mathcal O (\| a - a_j \|_{L^2 }).
\end{equation}
With the notation $ \langle b \rangle_S (z, \zeta) := \frac{1}{S} \int_{ 0}^S b (z + s \zeta) \ds $, $ b \in \CI (\TT^2) $, the last property in~\eqref{eq:propmu} and the fact that $ W_\infty $ is invariant under the flow shows that for any $ S > 0 $,
\[
\int_{W_\infty } a_j (z,\zeta) \dwmu_T (z,\zeta) =
\int_{W_\infty } \langle a_j \rangle_S (z, \zeta) \dwmu_T (z,\zeta).
\]
We note that
\begin{equation}\label{eq:Wm}
W_{m+1} \subset W_m, \ \ W_\infty := \bigcap_{ m =1}^\infty W_m = \{(z, \zeta) \tq |\zeta|=1,
\zeta \in \R^2 \setminus \mathbb Q^2 \}.
\end{equation}
For $ (z, \zeta) \in W_\infty $, unique ergodicity of the flow $ z \mapsto z + s \zeta $ shows that $ \langle a_j \rangle_S \to \langle a_j
\rangle := \int_{\TT^2} a_j (z) \dz / (2 \pi)^2. $ Fatou's Lemma then shows that
\[
\begin{split} \int_{W_\infty } a_j (z) \dwmu_T (z,\zeta)
& = \liminf_{ S \to \infty } \int_{W_\infty } \langle a_j \rangle_S (z, \zeta) \dwmu_T (z,\zeta) \\
& \geq \int_{W_\infty} \liminf_{ S \to \infty } \langle a_j \rangle_S (z, \zeta) \dwmu_T (z,\zeta) = \widetilde \mu_T (W_\infty) \langle a_j \rangle.
\end{split}
\]
Combining this with~\eqref{eq:muTaj} shows that
\[
\widetilde \mu_T (W_\infty) \leq \frac{ C \| a - a_j \|_{L^2 } } {\langle a_j \rangle} \rightarrow 0, \ \ j \to \infty,
\]
(since $ \| a \|_{L^2 } > 0 $ and $ a \geq 0 $, $ \langle a_j \rangle
\to \langle a \rangle > 0 $) which gives $ \widetilde \mu_T (W_\infty) = 0 $. But then~\eqref{eq:Wm} implies that $ \lim_{ m \to \infty } \widetilde \mu_T (W_m) = \widetilde \mu_T (W_\infty) = 0 $, concluding the proof.
\end{proof}

\subsection{Reduction to one dimension}

We start with the following

\begin{lemm}\label{l:RN}
Suppose that in~\eqref{eq:l2} $ m_0 $ is large enough so that $ \widetilde \mu_T (W_{m_0}) < T $ and that $ (z, \zeta_0) \in \supp (\widetilde \mu_T |_{ \complement W_{m_0} }) $. Then there exists $ F \in L^2 (\TT^2) $ such that
\begin{equation}\label{eq:muT0}
\widetilde \mu_T |_{ \TT^2 \times \{ \zeta_0 \} } = F \otimes \delta_{ \zeta = \zeta_0 }, \ \ \ \| F \|_{ L^2 (\TT^2) } \neq 0, \ \ F \geq 0.
\end{equation}
\end{lemm}

\begin{proof}
Let $ \pi : T^*\TT^2 \to \TT^2 $ be the natural projection map, $ \pi (x, \xi) = x $. Then, using~\eqref{eq:mT} and~\eqref{eq:defwm}, for any Lebesgue measurable set $ E \subset \TT^2 $,
\begin{equation}\label{eq:FmT}
\pi_* (\widetilde \mu_T |_{ \TT^2 \times \{ \zeta_0 \}}) (E)
\leq \pi_* (\widetilde \mu_T) (E) = \int_E m_T (z) \dz, \ \ \ m_T \in L^2.
\end{equation}
The Radon--Nikodym theorem then shows that $ \pi_* (\widetilde \mu_T |_{ \TT^2 \times \{ \zeta_0 \}}) = g m_T $ where $ g $ is measurable, $m_T$-a.e. finite. The inequality~\eqref{eq:FmT} gives $ F := g m_T \leq m_T $ almost everywhere which shows that $ F \in L^2 $.
\end{proof}

From now on we fix $ F $ and $ \zeta_0 $ such that~\eqref{eq:muT0} holds.

Using~\cite[Lemma~2.7]{BZ4} (see also~\cite[Figure~1]{BZ4}) we can assume (by changing the torus but not $ \Delta_z $) that $ \zeta_0 = (0, 1) $, $ z = (x, y) $, $ x \in \R/A_1 \ZZ $, $ y \in \R /B_1 \ZZ $, $ A_1/B_1 \in
\mathbb Q $. Abusing the notation we will keep the notation $ u_n $ and $
\mu $ for the transformed functions. Lemma~\ref{l:RN} and the invariance property in~\eqref{eq:propmu} (which implies that $ F (x, y) = g (x) $ is independent of $ y $) show now that
\begin{equation}\label{eq:muT00}
\begin{gathered}
(\widetilde \mu_T |_{ \TT^2 \times \{ (0,1) \} }) = g (x) \dx\dy \otimes \delta_{0} (\xi) \otimes\delta_1 (\eta), \ \ g \in L^2 (\TT^1), \ \ g \geq 0 \\
\int_{ \TT^2 } g (x) a (x, y) \dx \dy = 0, \ \
\| g \|_{ L^2 (\TT^2)} \neq 0.
\end{gathered}
\end{equation}
For $ m $ which will be chosen later (independently of the $ m_0 $ used in Lemma~\ref{l:RN}) we choose $ \chi = \chi_m \in \CIc (\R^2; [0,1])$ supported near $ | \zeta |= 1 $ and such that
\begin{equation}\label{eq:suppchi}
(\TT^2 \times \supp \chi) \cap \complement W_m = \{ (0, 1) \}, \ \ \ \chi(0,1) = 1.
\end{equation}
That is always possible as the set $ \{\zeta \tq (z, \zeta) \in \complement W_m \} $ is discrete (see the definitions in Lemma~\ref{l:2}).

We then define $ v_n := \chi (h D_z) u_n $ and $\nu := |\chi (\zeta)|^2 \mu \neq 0 $. Definition~\eqref{eq:mu} shows that $ \nu $ is the semiclassical defect measure associated to $ v_n (t) := e^{ i t \Delta } v_n = \chi (h D_z) e^{ it \Delta } u_n $ which in particular shows that
\begin{equation}\label{eq:beta}
\lim_{n\to \infty} \| v_n \|_{ L^2 (\TT^2) }^2 = T^{-1} \nu ([0, T] \times T^* \TT^2) = T^{-1} \int_{ \TT^2 } g (x) \dx \dy =: \beta > 0.
\end{equation}

The reduction to a one dimensional problem is based, as in~\cite{BZ4}, on a Fourier expansion in $ y $ (assuming $ B_1 = 2 \pi $ for notational simplicity):
\begin{equation}\label{eq:Four}
v_n (t) (x, y) := [e^{ i t \Delta } v_n] (x, y) =
\sum_{ k \in \ZZ } [e^{ i t \partial_x^2 } v_{n, k }] (x) e^{-it k^2 + i ky }
\end{equation}

We will now use a one dimensional analogue of Theorem~\ref{th.1} with $ W \in L^2 (\TT^1) $. It will be proved in Section~\ref{1D} below.
\begin{lemm}\label{l:3}
Suppose that $ b \in L^1 (\TT^1) $, $ b\geq 0 $, $\| b \|_{ L^1 } > 0 $ and that $ T > 0 $. Then there exists $ C $ such that for $ w \in L^2 (\TT^1) $,
\begin{equation}\label{eq:l3}
\| w \|_{L^2 (\TT^1) }^2 \leq C \int_0^T \int_{\TT^1} b (x) [e^{ i t \partial_x^2} w](x)|^2 \dx \dt.
\end{equation}
\end{lemm}

Let $ a_j $ be again given by~\eqref{eq:aj}. We apply~\eqref{eq:l3} to~\eqref{eq:Four} with $ b = \langle a \rangle_y := \frac{1}{ 2 \pi} \int_{ \TT^1} a (x, y) \dy $. This and~\eqref{eq:beta} give, for $ n $ large enough (and using Theorem~\ref{th.BBZ} to pass from $ a$ to $ a_j $ in the third line),
\begin{align*}
0 < \tfrac12 \beta & < \| v_n \|_{L^2 (\TT^2) }^2 = 2 \pi \sum_{k \in \ZZ } \| v_{n,k} \|_{ L^2 (\TT^1) }^2
\leq C' \int_0^T \int_{\TT^1} \langle a \rangle_y (x) \sum_{ k \in \ZZ} | [e^{ it \partial_x^2 } v_{n,k}] (x) |^2 \dx \dt
\\
& = C \int_0^T \int_{ \TT^2 } \langle a \rangle_y (x) |[e^{ it \Delta } v_n] (x, y) |^2 \dx \dy \dt \\
& = C \int_0^T \int_{\TT^2} \langle a_j \rangle_y (x) | [e^{ it \Delta } v_n] (x, y) |^2 \dx \dy \dt +
\mathcal O (\| a - a_j \|_{L^2 (
\TT^2) }) \\
& \longrightarrow C \int_{T^* \TT^2 }
\langle a_j \rangle_y (x) \dwnu_T (x, y, \xi, \eta) + \mathcal O (\| a - a_j \|_{L^2 (
\TT^2) }), \ \ n \longrightarrow \infty,
\end{align*}
where $ \widetilde \nu_T = |\chi (\xi, \eta)|^2 \widetilde \mu_T $ (see~\eqref{eq:defwm}). In particular for every $ j $,
\begin{equation}\label{eq:ajn}
0 < \alpha \leq \int_{T^* \TT^2 }
\langle a_j \rangle_y (x) \dwnu_T (x, y, \xi, \eta) + \mathcal O (\| a - a_j \|_{L^2 (
\TT^2) }), \ \
\alpha := \frac{ \beta}{ C }.
\end{equation}
We now decompose the integral in~\eqref{eq:ajn} as $ I_1 + I_2 $ and use~\eqref{eq:muT00}:
\begin{equation}\label{eq:aj1}
\begin{split}
I_1 := \int_{ \TT^2 \times \{ (0, 1) \} } \langle a_j \rangle_y (x) \dwnu_T (x, y, \xi, \eta) & =
\int_{ \TT^2 } g (x) a_j (x, y) \dx \dy \\
& =
\int_{ \TT^2 } g (x) (a_j (x, y) - a (x, y)) \dx \dy \\
& \leq {\sqrt{ 2 \pi }} \| g \|_{ L^2 (\TT^1) } \| a_j - a \|_{L^2 (\TT^2) }.
\end{split}
\end{equation}
We then use~\eqref{eq:suppchi} to estimate the remainder:
\[
\begin{split}
I_2 & := \int_{ (\xi, \eta) \neq (0, 1) }
\langle a_j \rangle_y (x) \dwnu_T (x, y, \xi, \eta)
\leq \int_{W_m } \langle a_j \rangle_y (x) \dwnu_T (x, y, \xi, \eta) \\
&\, \leq \| a_j \|_{ L^\infty } \, \widetilde \nu_T (W_m).
\end{split}
\]
We now combine these two estimates with~\eqref{eq:ajn} to obtain:
\[
0 < \alpha \leq K \| a_j - a \|_{L^2 (\TT^2) } + \| a_j \|_{ L^\infty } \, \widetilde \nu_T (W_m),
\]
where the constant $ K $ does not depends on $ \chi $ and $ m $.

Hence, we first choose $ j $ large enough so that $ K \| a_j - a \|_{L^2 (\TT^2) } < \alpha/2 $ and then $ m $ large enough and $ \chi $ satisfying~\eqref{eq:suppchi} so that~\eqref{eq:l2} gives
\[
\| a_j \|_{ L^\infty } \, \widetilde \nu_T (W_m) <
\epsilon \| a_j \|_{ L^\infty } < \alpha/2.
\]
This provides a contradiction and proves Proposition~\ref{p:semi}.

\subsection{One dimensional estimate}\label{1D}

We now prove Lemma~\ref{l:3}. The semiclassical part proceeds along the lines of the proof of Proposition~\ref{p:semi}. The derivation of~\eqref{eq:l3} from the semiclassical estimate follows the same arguments needed in Section~\ref{obs} and we will refer to that section for details.

\begin{proof}[Proof of Lemma~\ref{l:3}]
We start with a semiclassical statement: for every $ T $ there exist $ K$, $ \rho_0 $ and $ h_0 $ such that for $ 0 < h < h_0 $ and $ 0 < \rho < \rho_0 $ we have the analogue of~\eqref{eq:semi}:
\begin{multline}\label{eq:semi1}
\| \pi_{h, \rho} u_0 \|_{ L^2 (\TT^1) }^2 \leq K \int_0^T \!\!\int_{ \TT^1}b (z) | e^{ i t \Delta} \pi_{h, \rho} u_0 (z) |^2 \dz \dt,
\\
\pi_{h, \rho} (u_0) := \chi \left(\frac{ h^2 D_x^2 - 1} {\rho} \right) u_0 \,.
\end{multline}
We proceed by contradiction which leads to an analogue of~\eqref{eq:contr} and then to a measure $ \omega_T $ analogous to $ \widetilde \mu_T $ (see~\eqref{eq:defwm}) on $ T^* \TT^1 $ and satisfying: $ \supp \omega_T \subset \{ \xi = \pm 1 \} $, $\partial_x \omega_T = 0$, where the derivative is taken in the distributional sense.

From~\cite[Proposition~2.1]{BBZ}\footnote{See \url{https://math.berkeley.edu/~zworski/corr_bbz.pdf} for a corrected version. That correction is relevant only when treating the irrational case: see~\cite{BZ4} and~\cite{BBZ}. That proposition is an easy one dimensional analogue of Theorem~\ref{th.BBZ}.} and the argument in Lemma~\ref{l:RN} (with weak convergence in $L^2 $ replaced by the weak$*$ convergence in $ L^\infty = (L^1)^* $) we obtain
\[
d \omega_T = \sum_{ \pm } f_\pm (x) \dx \otimes \delta_{\pm 1} (\xi) \dxi, \ \ \ f_\pm \in L^\infty (\TT^1), \ \ f_\pm \geq 0.
\]
But the fact that $ \partial_x \omega_T = 0 $ and the analogue of~\eqref{eq:l1} show that
\[
f_\pm (x) = c_\pm \geq 0, \ \ c_+ + c_- > 0, \ \ (c_+ + c_-)\int _{\TT^1 } b (x) \dx = 0,
\]
which is a contradiction proving~\eqref{eq:semi1}.

From the semiclassical estimate we obtain
\[
\| u_0\|_{L^2 (\TT^1) } \leq C
\int_0^T \!\!\int_{ \TT^1}b (z) | e^{ i t \Delta} u_0 (z) |^2 \dz \dt + C \| u_0 \|_{ H^{-2}(\TT^1)}.
\]
That is done by the same argument recalled in Section~\ref{dya} below. Finally the error term $ \| u_0 \|_{ H^{-1}(\TT^1) } $ is removed (see Section~\ref{eli} for review of the procedure for doing (applying~\cite[Proposition~2.1]{BBZ} again)).
\end{proof}


\section{Observability estimate}\label{obs}

To prove Theorem~\ref{th.1} we first prove a weaker statement involving an error term:

\begin{prop}\label{p:weak}
Suppose that $ W \in L^4 (\TT^2) $, $ a \geq 0 $ and $ \| W \|_{ L^4 } \neq 0 $. Then for any $ T > 0 $, there exists $ K $ such that for $ u \in L^2 $,
\begin{equation}\label{eq:weak}
\| u \|_{ L^2 (\TT) }^2 \leq C \int_0^T \int_{\TT^2 } | W(z) e^{ i t \Delta } u (z) |^2 \dz \dt + C \| u \|^2_{ H^{-2} (\TT^2) }.
\end{equation}
\end{prop}

\subsection{Dyadic decomposition}\label{dya}
The proof of~\eqref{eq:weak} uses a dyadic decomposition as in~\cite[Section~5.1]{BBZ} and~\cite[Section~4]{BZ4} and we recall the argument adapted to the setting of this paper. For that let $ 1= \varphi_{0} (r)^2 + \sum_{k=1}^{\infty} \varphi_{k} (r)^2$, where for $ k \geq 1 $
\begin{gather*}
\varphi_{k} (r) := \varphi(R ^{-k} |r|), \enspace R > 1, \\ \varphi \in \CIc ((R^{-1}, R) ; [0, 1]), \quad (R^{-1}, R) \subset \{ r \tq \chi(r / \rho) \geq
\textstyle{\frac12} \},
\end{gather*}
with $ \chi $ and $ \rho $ same as in~\eqref{eq:Pih} and~\eqref{eq:semi}. Then, we decompose $u_0$ dyadically: $\|u_0 \|^2_{L^2} = \sum_{k=0}^\infty \|\varphi_{k}(- \Delta) u_0 \|_{L^2}^2, $ which will allow an application of Proposition~\ref{p:semi}.

\begin{proof}[Proof of Proposition~\ref{p:weak}]
Let $\psi \in \CIc ((0, T) ; [0, 1]) $ satisfy $\psi (t) > 1/2$, on $ T/3 < t < 2T / 3 $. Proposition~\ref{p:semi} applied with $ a = W^2 $, $ u_0 = e^{ i (T/3) \Delta } \Pi_{h, \rho } u_0 $ and with $ T $ replaced by $ T/3 $, shows that
\begin{equation}\label{eq:piece}
\|\Pi_{h, \rho } u_0\|_{L^2}^2 \leq K \int_\R \psi(t)^2 \| W \mathrm{e}^{\mathrm{i} t \Delta } \Pi_{h, \rho } u_0\|_{L^2 (\TT^2) }^2 \dt, \quad 0 < h < h_0.
\end{equation}
Taking $K $ large enough so that $R^{-K } \leq h_{0}$ we apply~\eqref{eq:piece} to the dyadic pieces:
\begin{align*}
\|u_0 \|^2_{L^2} & =\sum_{ k \in \ZZ } \|\varphi_{k}(-\Delta)u_0\|_{L^2}^2 \\
&\leq \sum_{k=0}^{K } \|\varphi_{k}(-\Delta)u_0\|_{L^2}^2 +C
\sum_{k=K +1}^\infty \int_{0}^T \psi (t) ^2 \| W \varphi_{k}(-\Delta)\,
\mathrm{e}^{\mathrm{i}t \Delta } u_0\|_{L^2(\TT^2) } ^2 \dt \\
& = \sum_{k=0}^{K } \|\varphi_{k}(-\Delta)u_0\|_{L^2}^2 + C
\sum_{k=K +1}^\infty \int_\R \| \psi (t) W \varphi_{k}(D_t)
\,\mathrm{e}^{\mathrm{i}t \Delta } u_0\|_{L^2(\TT^2) } ^2 \dt.
\end{align*}
In the last equality we used the equation and replaced $ \varphi (-\Delta) $ by $
\varphi (D_t) $.

We need to consider the commutator of $ \psi \in \CIc ((0, T)) $ and $
\varphi_{k} (D_t)\,{=}\,\varphi (R^{-k} D_t) $. If $
\widetilde \psi \in \CIc ((0,T))$ is equal to $ 1 $ on $
\supp \psi$ then the semiclassical pseudo-differential calculus with $ h = R ^{-k}$ (see for instance~\cite[Chapter~4]{EZB}) gives
\begin{equation}\label{eq:wideps}
\psi (t) \varphi_{k} (D_t) = \psi (t) \varphi_{k} (D_t) \widetilde \psi (t) + E _k (t, D_t),\quad
\partial^\alpha E_k = {\mathcal O} (\langle t \rangle^{-N } \langle
\tau \rangle^{-N} R ^{- N k }),
\end{equation}
for all $ N$ and uniformly in $ k $.

The errors obtained from $ E_k $ can be absorbed into the $
\| u_0 \|_{ H^{-2} (\TT^2) } $ term on the right-hand side. Hence we obtain,
\begin{align*}
\|u_0 \|^2_{L^2} &\leq C \|u_0 \|_{H^{-2} (\TT^2)} ^2+ C \sum_{k=0}^\infty \int_{0}^T
\| \psi (t) \varphi_{k}(D_t)\, \widetilde \psi (t) W \mathrm{e}^{ \mathrm{i}t \Delta }u_0 \|^2_{L^2(\TT^2)} \dt\\
& \leq\widetilde C \|u_0 \|_{H^{-2} (\TT^2)} ^2+ K \sum_{k=0}^\infty
\langle \varphi_{k}(D_t)^2 \widetilde \psi (t) \, W \mathrm{e}^{\mathrm{i}t \Delta }u_0, \psi(t)^2 \widetilde \psi(t)\, W \mathrm{e}^{\mathrm{i}t \Delta}u_0, \rangle_{L^2 (\R_t \times \TT^2)} \\
& \leq \widetilde C \|u_0 \|_{H^{-2} (\TT^2)} ^2 + K \int_\R \| \widetilde \psi (t) \, W \mathrm{e}^{\mathrm{i}t \Delta} u_0\|^2_{L^2(\TT^2)} \dt \\
& \leq \widetilde C \|u_0 \|_{H^{-2} (\TT^2)} ^2 + K \int_{0}^T \|\, W\mathrm{e}^{\mathrm{i}t \Delta } u_0\, \|^2_{L^2(\TT^2)} \dt,
\end{align*}
where the last inequality is~\eqref{eq:weak} in the statement of the proposition.
\end{proof}

\subsection{Elimination of the error term}\label{eli}

We now eliminate the error term on the right hand side of~\eqref{eq:weak}. For that we adapt the now standard method of Bardos--Lebeau--Rauch~\cite{BLR} just as we did at the end of~\cite[Section~4]{BZ4}. The argument recalled there shows that if
\begin{equation}\label{eq:defN0}
N := \{ u \in L^2 (\TT^2) \tq W e^{ i t \Delta } u \equiv 0
\ \text{ on $ (0, T) \times \TT^2 $}\}
\end{equation}
is non-trivial then since $ iW e^{it\Delta } \Delta u= \partial_t W e^{ i t \Delta } u \equiv 0$ on $ (0, T) \times \TT^2 $, then $N$ is invariant by the action of $\Delta$, and hence it contains a nontrivial $ w \in L^2 (\TT^2) $ such that for some $ \lambda $,
\[
(- \Delta - \lambda) w = 0, \ \ \ W w \equiv 0.
\]
But then $ w $ is a trigonometric polynomial vanishing on a set of positive measure which implies that $ w \equiv 0 $. Hence
\begin{equation}\label{eq:N0}
N = \{ 0 \}.
\end{equation}

\begin{proof}[Proof of Theorem~\ref{th.1}]
Suppose the conclusion~\eqref{eq:th1} were not to valid. Then there exists a sequence $ u_n \in L^2 (\TT^2) $ such that
\begin{equation}\label{eq:unW}
\| u_n \|_{ L^2 (\TT^2) } = 1, \ \ \
\| W e^{it \Delta } u_n \|_{ L^2 ((0, T) \times \TT^2) }\to 0, \ \ n
\to \infty.
\end{equation}
By passing to a subsequence we can then assume that $ u_n $ converging weakly in $ L^2 (\TT^2) $ and strongly in $ H^{-2} (\TT^2) $ to some $ u \in L^2 $. From Proposition~\ref{p:weak} we would also have
\[
1 = \| u_n \|_{ L^2 (\TT^2) }^2 \leq C
\int_0^T \| W e^{it \Delta } u_n \|_{ L^2 (\TT^2) }^2 \dt + C \| u_n \|^2_ {H^{-2} (\TT^2) }.
\]
Hence,
\begin{equation}\label{eq:cont1}
1 \leq C \lim_{n \to \infty } \| u_n \|_{H^{-2} (\TT^2) } = C\| u \|_{ H^{-2} (\TT^2) }
\ \ \Longrightarrow \ \ u \not \equiv 0.
\end{equation}
Let $ W_j \in \CI (\TT^2) $ satisfy $ \| W - W_j \|_{L^4 (\TT^2) }
\to 0 $. For $ \varphi \in \CIc ((0, T) \times \TT^2) $, due to distributional convergence, Theorem~\ref{th.BBZ} and~\eqref{eq:unW},
\[
\begin{split} |\langle W_j e^{ it\Delta } u, \varphi \rangle| & = \lim_{ n \to \infty }
|\langle e^{ it \Delta } u_n, W_j \varphi \rangle| \\
&\leq
\lim_{ n \to \infty }
\left(| \langle (W_j - W) e^{ i t \Delta } u_n, \varphi \rangle | +
\langle W e^{ it \Delta } u_n, \varphi \rangle \right) \\ & \leq
\|\varphi \|_{ L^2 } \| (W_j - W) e^{it \Delta } u_n \|_{L^2 ((0, T) \times
\TT^2) }
\leq C\|\varphi \|_{L^2 } \| W_j - W \|_{L^4(\TT^2) }.
\end{split}
\]
On the other hand the same argument shows that
\[
| \langle W e^{it \Delta }u, \varphi \rangle | \leq | \langle W_j e^{ it \Delta } u, \varphi \rangle | + C \| \varphi\|_{L^2} \| W_j - W \|_{L^4(\TT^2) }.
\]
Combining the two inequalities we see that 
\[
|\langle W e^{it \Delta }u, \varphi \rangle | \leq C
\lim_{ j \to \infty } \| W_j - W \|_{L^4(\TT^2) } = 0
\]
which means that $ W e^{ it \Delta } u \equiv 0 $. Thus $ u \in N $ given by~\eqref{eq:defN0} and by~\eqref{eq:N0}, $ u = 0 $. This contradicts~\eqref{eq:cont1} completing the proof.
\end{proof}


\section{The HUM method: proofs of Theorems~\ref{th.2} and~\ref{th.3}}\label{HUM}

We now show the equivalence of the stabilization, control and observability properties in our context. The proof is a variation on the classical HUM method~\cite{Li}, but since our damping and localization functions are \emph{not} in $L^\infty$ it requires additional care.
\begin{prop}\label{p:hum}
The following are equivalent (for fixed $T>0$).
\begin{enumerate}
\item \label{i}
Let $a \in L^2(\TT^2;\R), \| a \|_{L^2} >0$. For any $u_0\in L^2(\TT^2)$ there exists $f \in\linebreak L^4(\TT^2 ; L^2(0,T))$ such that the solution $u$ of~\eqref{eq:con} satisfies $u|_{t=T} =0$
\item \label{ii}
Let $a \in L^4(\TT^2;\R), \| a \|_{L^4} >0$. For any $u_0\in L^2(\TT^2)$ there exists $f \in\linebreak L^2((0,T)\times \TT^2)$ such that the solution $u$ of~\eqref{eq:con} satisfies $u |_{t=T} =0$
\item \label{iii}
Let $a \in L^4(\TT^2;\R), \| a \|_{L^4} >0$. Then there exists $C>0$ such that for any $v_0 \in L^2(\TT^2)$,
\begin{equation}\label{observation}
\| v_0 \|_{L^2(\TT^2)}^2 \leq C \| a e^{it \Delta} v_0 \| _{L^2((0,T) \times \TT^2)}.
\end{equation}
\end{enumerate}
\end{prop}

\begin{proof}
We first prove that~\eqref{ii} and~\eqref{iii} are equivalent, and for that we follow the HUM method. Define the map
\[
R: f \in L^2((0,T) \times \TT^2) \mapsto R f = u |_{t=0},
\]
where $u$ is the solution of the final value problem
\[
(i \partial_t + \Delta) u (z) = a(x) 1_{(0,T)}f \in L^{4/3}(\TT^2; L^2(0,T)), \quad u |_{T=0} =0.
\]
By Theorem~\ref{th.BBZ} $ R : L^{2}(\TT^2; L^2(0,T)) \to L^2 (\TT^2) $ and
\begin{equation}\label{eq:2}
\eqref{ii}\Longleftrightarrow R (L^{2}(\TT^2; L^2(0,T))) = L^2 (\TT^2).
\end{equation}

Again by Theorem~\ref{th.BBZ}, $e^{it \Delta} v_0 \in L^4(\TT^2; L^2(0,T)) $ for $ v_0 \in L^2 (\TT^2) $, we define
\[
S: v_0 \in L^2(\TT^2) \mapsto 1_{(0,T)} \times a e^{it \Delta} v_0 \in L^2((0,T) \times \TT^2),
\]
and
\begin{equation}\label{eq:3}
\eqref{iii}\Longleftrightarrow \ \exists K \, \forall v_0 \in L^2 (\TT^2) \ \
\| v_0 \|_{L^2 (\TT^2) } \leq K\| S v_0 \|_{L^2 ((0, T) \times \TT^2)}.
\end{equation}
To relate $ R $ and $ S $ we integrate by parts:
\[
\begin{split}
\int_0^T \!\!\int_{\TT^2} af \overline {v} \dx \dt &= \int_0^T \!\!\int_{\TT^2} (i \partial_t + \Delta) u \overline {v} \dx \dt = i \left[\int_{\TT^2} u \overline{v} \dx \right] _0 ^T +\int_0^T \!\!\int_{\TT^2} u \overline {(i \partial_t + \Delta)v} \dx \dt
\\ & = -i \int_{\TT^2} u \overline{v} \dx|_{t=0},
\end{split}
\]
which is the same as
\begin{equation}\label{eq:SR}
\bigl(f, Sv_0\bigr)_{L^2((0,T)\times \TT^2)} = -i \bigl(Rf, v_0\bigr)_{L^2(\TT^2)}.
\end{equation}

Let us assume~\eqref{ii}. By~\eqref{eq:2} and the closed graph theorem there exists $\eta>0$ such that the image of the unit ball in $L^2((0,T) \times \TT^2)$ by $R$ contains the ball $ \{ v_0 \in L^2 (\TT^2) \tq \| v_0 \|_{L^2} \leq \eta \} $. Hence for all $ v_0 \in L^2 (\TT^2) $ there exists $ f \in L^2((0,T) \times \TT^2)$, such that
\[
\| f \|_{ L^2 ((0,T) \times \TT^2) } \leq \frac 1 \eta \| v_0\|_{ L^2}, \ \ \ R f = v_0.
\]
Hence, using~\eqref{eq:SR},
\begin{equation}
\begin{split}
\| v_0 \| _{L^2(\TT^2)}^2 & = i \bigl(f, Sv_0\bigr)_{L^2((0,T)\times \TT^2)} \leq \| f\|_{L^2((0,T)\times \TT^2} \| Sv_0 \|_{L^2((0,T)\times \TT^2)}
\\
& \leq \frac 1 \eta \| Sv_0 \|_{L^2((0,T)\times \TT^2)}
\| v_0 \| _{L^2(\TT^2)},
\end{split}
\end{equation}
and by~\eqref{eq:3}, {\eqref{iii}} follows.

On the other, assume that {\eqref{iii}} holds. By~\eqref{eq:3}, the operator
\[
-i R \circ S : L^2(\TT^2) \rightarrow L^2(\TT^2)
\]
is continuous and,
by~\eqref{eq:SR}, there exists $ C > 0 $ such that for all $ v_0 \in L^2(\TT^2) $,
\[
\bigl(-i R \circ S v_0, v_0 \bigr) _{L^2(\TT^2)} = \bigl(Sv_0, S v_0\bigr)_{L^2((0,1) \times \TT^2) } \geq \frac 1 C \| v_0 \|_{L^2(\TT^2)}^2.
\]
Consequently $ -i R \circ S $ is an injective bounded self-adjoint operator, hence bijective. This in turn shows that $R$ is surjective and in view of~\eqref{eq:2} proves~\eqref{ii}.

We conclude by showing that~\eqref{i} is equivalent to~\eqref{ii}. Since $ L^4 (\TT^2, L^2 (0, T)) \subset L^2 ((0, T) \times
\TT^2)$, \eqref{ii} is immediate from~\eqref{i}.

On the other hand, suppose that $ a \in L^2 (\TT^2) $ so that $ |a|^{\frac12} \in L^4 (\TT^2) $. We can then apply~\eqref{ii} with $ a $ replaced by $ |a|^{\frac12 } $. The proof above shows that we can take $ \tilde f = Sv_0 = |a|^{\frac12} e^{it \Delta} v_0$, $ v_0 = (R \circ S)^{-1} u_0 $ to obtain a solution to $ (i \partial_t + \Delta) u = |a|^{\frac12} \tilde f $, $ u|_{t=0} = u_0 $ with $ u |_{t=T}=0 $. Theorem~\ref{th.BBZ} shows that $ f := e^{ i \arg(a)} e^{it \Delta} v_0 \in L^4(\TT^2; L^2(0,T)) $ and that~\eqref{eq:con} holds with $ u |_{t=T}=0 $. Since now $ a \in L^2 (\TT^2) $ that is~\eqref{i}.
\end{proof}

In view of Theorem~\ref{th.1} this proves Theorem~\ref{th.2} and provides some additional versions of it. We now turn to the damped Schrödinger equation.

\begin{proof}[Proof of Theorem~\ref{th.3}]
For $a \in L^2$ put $H := (- i \Delta +a)$. Then for $ \lambda \in \CC $, $ H + \lambda : H^2 (\TT^2) \to L^2 (\TT^2) $ is a Fredholm operator of index zero. In fact,
\[
H + \lambda = (I + b (i \Delta + 1)^{-1}) (i \Delta + 1), \ \ b:= a + \lambda - 1.
\]
Since $ (i \Delta + 1) :H^2 (\TT^2) \to L^2 (\TT^2) $ is an isomorphism, this follows from the fact that for $ b \in L^2 (\TT^2) $, $ b (i \Delta + 1)^{-1} : L^2 \to L^2 $ is compact. This can be seen as follows: if $ b_j \in C^\infty (\TT^2) $ and $ \| b_j - b\|_{L^2 } \to 0 $ then
\[
\begin{split} \| (b - b_j) (i \Delta +1)^{-1} \|_{ L^2 \to L^2 } & \leq
\| b - b_j \|_{ L^2 } \| (i \Delta + 1)^{-1} \|_{ L^2 \to L^\infty } \\
&
\leq C \| b - b_j \|_{ L^2 } \| (i \Delta + 1)^{-1} \|_{ L^2 \to H^2 }
\to 0,
\end{split}
\]
and $ b_j (i \Delta + 1)^{-1} : L^2 (\TT^2) \to H^2 (\TT^2) $ is compact as an operator $ L^2 (\TT^2)\to L^2 (\TT^2) $.


If $ a \geq 0 $ and $ \lambda \in \R $ we also have.
\[
\Re \bigl((Hu + \lambda) u, u \bigr)_{L^2} =
\lambda \| u \|_{L^2(\TT^2)} ^2 +
\int_{\TT^2} a |u|^2 (x) \dx \geq \lambda \| u \|_{L^2(\TT^2)} ^2, \ \ u \in H^2 (\TT^2).\ \
\]
The Fredholm property shows that for $\lambda >0$ the equation $(H + \lambda) u =f \in L^2(\TT^2)$ can be solved with $ \| u \|_{ L^2 } \leq \lambda^{-1} \| f \|_{L^2 } $. The Hille--Yosida theorem then shows that $ H $ defines a strongly continuous semigroup $ [0, \infty) \ni t \mapsto \exp (- t H) $. Furthermore, when $u_0 \in H^2$,
\[
u(t) := \exp (- t H) u_0 \in C^1([0, \infty) ; L^2(\TT^2))\cap C^0([0, \infty) ; H^2(\TT^2)).
\]
We then check that
\begin{equation}\label{damping}
\begin{gathered}
\| u(t) \|^2_{L^2(\TT^2)} = \| u_0 \|^2_{L^2(\TT^2)}- \int_0 ^t \int_{\TT^2} a(x) |u|^2(s,x) \dx \ds \leq \| u_0 \|^2_{L^2(\TT^2)},\\
u(t) = e^{it \Delta} u_0 - \int_0^t e^{i(t-s)\Delta} (au)(s) \ds.
\end{gathered}
\end{equation}
Let $a_j \in C^0 (\TT^2) $ and $ \| a_j - a\|_{L^2 (\TT^2) } \to 0 $, $ j \to \infty $. Using the second expression in~\eqref{damping} and Theorem~\ref{th.BBZ} we obtain
\[
\begin{split}
\| u\|_{L^4(\TT^2; L^2(0,T))} & \leq C \| u_0 \|_{L^2(\TT^2)} + C\| a_j u \|_{L^1((0,T); L^2(\TT^2))} + C\| (a- a_j)u \|_{L^{4/3}(\TT^2; L^2(0,T))}\\
& \leq C \| u_0 \|_{L^2(\TT^2)} +C T\| a_j\|_{L^\infty} \| u\|_{L^{\infty}((0,T); L^2(\TT^2))}
\\ & \qquad +C\| a-a_j\|_{L^2(\TT^2)} \| u \|_{L^{4}(\TT^2; L^2(0,T))}
\end{split}
\]
Taking $ j$ large enough so that $C\| a-a_j\|_{L^2(\TT^2)} \leq \frac12$ and recalling the first line of~\eqref{damping} we get
\[
\| u\|_{L^4(\TT^2; L^2(0,T))} \leq C' \| u_0 \|_{L^2(\TT^2)}, \ \
u_0 \in H^2 (\TT^2).
\]
Since $ H^2 $ is dense in $L ^2 $, this remains true for initial data $u_0\in L^2(\TT^2)$ and consequently, for $ a \in L^2 $, we get that
\begin{equation}\label{esti-ap}
\int_0 ^t \int_{\TT^2} a(x) |u|^2(s,x) \dx \ds\leq C \| u_0 \|_{L^2(\TT^2)}.
\end{equation}
By simple integration by parts~\eqref{damping} remains valid for $u_0\in H^2$, and consequently from~\eqref{esti-ap} it remains true for $u_0\in L^2$. Now, if for some $ T > 0 $,
\begin{equation}\label{eq:u00}
\| u_0 \|_{L^2(\TT^2)} \leq C \int_0 ^T \int_{\TT^2} a(x) |u|^2(t,x) \dx \dt,
\end{equation}
where $u$ is the solution of~\eqref{damped}, then~\eqref{damping} and semigroup property show that $ \| u (k T) \|_{ L^2 (\TT^2)}^2 \leq (1 - 1/C)^N \| u_0 \|_{ L^2 (\TT^2) }^2$, and the exponential decay~\eqref{dampedbis} follows.

For any fixed $T>0$~\eqref{eq:u00} is the same as~\eqref{eq:th1} with $ W = a^{\frac12} $, except that here $u$ is the solution of the damped Schrödinger equation, while in~\eqref{eq:th1} it is the solution of the free Schrödinger equation.

We now claim that~\eqref{eq:th1}, $ W = a^{\frac12} $, implies~\eqref{eq:u00}. In fact, suppose that~\eqref{eq:u00} is not true. Then, there exists a sequence $ u_{0,n} \in L^2 $,
\[
\| u_{0,n} \|_{L^2} =1, \ \ \ (i \partial_t + \Delta)u_n = a u_n, \ \ u_n|_{t=0} = u_{0,n}, \ \
\| a^{\frac12} u_n \|_{ L^2 ((0, T) \times \TT^2) } \to 0, \ n\to \infty.
\]
Then
\[
\| a u_n \|_{ L^{\frac43} (\TT^2; L^2 ((0, T))) } =
\| a^{\frac12} a^{\frac12} u \|_{ L^{\frac43} (\TT^2; L^2 ((0, T))) }
\leq \| a^{\frac12} \|_{L^4(\TT^2) }
\| a^{\frac12} u_n \|_{ L^2 ((0, T) \times \TT^2) } \to 0,
\]
and Theorem~\ref{th.BBZ} shows that $ u_n = e^{ it \Delta} u_{0,n} + e_n$, $ \| e_n \|_{ L^4 (\TT^2 ; L^2 ((0, T))) } \to 0 $. But then, using~\eqref{eq:th1},
\[
\begin{split} 0 & = \limsup_{n \to \infty } \| a^{\frac12} u_n \|_{L^2 ((0, T) \times \TT^2) }
\\ &
\geq \limsup_{n \to \infty } \left(\| a^{\frac12} e^{ it \Delta } u_{ 0,n} \|_{ L^2 ((0, T) \times \TT^2) } - \| a^{\frac12} \|_{L^4 }
\| e_n \|_{ L^4 (\TT^2 ; L^2 ((0, T))) } \right) \\
& \geq
\limsup_{n \to \infty } \| a^{\frac12} e^{ it \Delta } u_{ 0,n} \|_{ L^2 ((0, T) \times \TT^2) } \geq c
\limsup_{n\to \infty } \| u_{0,n}\|_{L^2 (\TT)} = c > 0
\end{split}
\]
which gives a contradiction. Hence~\eqref{eq:u00} holds and that completes the proof.
\end{proof}

\appendix
\section*{Appendix}
\refstepcounter{section}

To see that Theorem~\ref{th.1} for $ T > \pi $ and rational tori follows from~\cite[Theorem~1.2]{Jak} assume that $ \TT^2 = (\R/2\pi \ZZ)^2 $. We then write $ u (z) = \sum_{ \lambda } u_\lambda $, where the sum of is over distinct eigenvalues of $ - \Delta $ (and $u_\lambda$ is the projection of $u$ on the corresponding eigenspace). By Ingham's inequality~\cite{ing} (this is where $ T > \pi $ is used),
\[
\int_0^T \| W e^{ it \Delta} u \|^2_{ L^2 (\TT^2)} =
\int_{\TT^2 } \int_0^T \left| \sum_{ \lambda \in \NN } W(z) u_\lambda (z) e^{i t \lambda } \right|^2 \dt \dz \geq B
\int_{\TT^2} \sum_{ \lambda \in \NN } | W (z) u_\lambda (z) |^2 \dz.
\]
Hence, \eqref{eq:th1} follows from the estimate,
\begin{equation}\label{eq:Jak}
\sum_{ \lambda } \| u_\lambda\|^2_{L^2 (\TT^2) } \leq C \int_{\TT^2} \sum_{ \lambda \in \NN } | W (z) u_\lambda (z) |^2 \dz,
\end{equation}
which it turn follows from a pointwise estimate:
\begin{equation}\label{eq:Jak1}
\| u_\lambda\|^2_{L^2 (\TT^2) } \leq C \int_{\TT^2} | W (z) u_\lambda (z) |^2 \dz, \ \ \ - \Delta u_\lambda = \lambda u_\lambda.
\end{equation}

\begin{proof}[Proof of~\eqref{eq:Jak1}]
We start with the observation that the zero set of a non-trivial trigonometric polynomial $ p (z) $ has measure zero and hence,
\begin{equation}\label{eq:trig}
\int_{\TT^2} | W (z) p(z) |^2 \dz > 0.
\end{equation}
In particular that holds for any fixed eigenfunction of $ - \Delta $.

To prove~\eqref{eq:Jak1} we proceed by contradiction, that is we assume that there exists a sequence of $ e_n $'s, such that
\begin{equation}\label{eq:en}
\|e_n \|_{L^2}^2 = 1, \ \ \ \| W e_n \|_{L^2}^2 \rightarrow 0, \ \ - \Delta e_n =
\lambda_n e_n.
\end{equation}

Suppose first that $ \lambda_n $ are bounded. We can then assume that $ \lambda_n \to \lambda $. From~\eqref{eq:en} we see that $ e_n $ are bounded in $ H^2 $ and hence we can assume that $ e_n \to e $ in $ H^1 $ and, as $ H^1 \subset L^4 $, also in $ L^4 $. Then~\eqref{eq:en} shows that $ -\Delta e = \lambda e$, $ \| e\|_{L^2 } = 1 $, $ \| W e \|_{L^2 } = 0$, which contradicts~\eqref{eq:trig}.

Hence we can assume (by extracting a subsequence) that $ \lambda_n \to \infty $ in~\eqref{eq:en}. We can then assume that the sequence of probability measures $ |e_n|^2 \dx $ converges weakly to a measure $\nu$. According to~\cite[Theorem~1.2]{Jak}, $ \nu = p (z) \dz $ where $ p$ is a non-negative trigonometric polynomial, $ \int p (z) \dz = 1 $.

Let $ f_k\in C^0$, $ f_k \geq 0$, converge to $|W|^2 $ in $L^2$. From Zygmund's bound on the $L^4$ norm of $e_n$~\eqref{eq:zyg0}, we get
\[
\limsup_{n\rightarrow + \infty} | \int (f_k - |W |^2)|e_n|^2 (x) \dx | \leq C \| f_k - |W|^2\|_{L^2},
\]
and from the weak convergence $ \lim_{n\rightarrow + \infty} \int f_k|e_n|^2 (x) \dx = \int f_k (x) (x)\dx.$ We deduce
\[
0= \lim_{n\rightarrow + \infty} \int |W e_n|^2 (x) \dx = \int |W (x)|^2 (x)\dx.
\]
This again contradicts~\eqref{eq:trig}.
\end{proof}

\begin{thebibliography}{MMMMMM}
\bibitem[AFKM15]{AFM} Anantharaman, Nalini; Fermanian-Kammerer, Clotilde; Macià, Fabricio Semiclassical completely integrable systems: long-time dynamics and observability via two-microlocal Wigner measures, Am. J. Math., Volume 137 (2015) no. 3, pp. 577-638 \href{https://ahl.centre-mersenne.org/item/AHL_2019__2__331_0/#AFM}{Publisher reference}.
\bibitem[AJM12]{AJM} Aïssiou, Tayeb; Jakobson, Dmitry; Macià, Fabricio Uniform estimates for the solutions of the Schrödinger equation on the torus and regularity of semiclassical measures, Math. Res. Lett., Volume 19 (2012) no. 3, pp. 589-599 \href{https://ahl.centre-mersenne.org/item/AHL_2019__2__331_0/#AJM}{Publisher reference}.
\bibitem[AM14]{AM} Anantharaman, Nalini; Macià, Fabricio Semiclassical measures for the Schrödinger equation on the torus, J. Eur. Math. Soc., Volume 16 (2014) no. 6, pp. 1253-1288 \href{https://ahl.centre-mersenne.org/item/AHL_2019__2__331_0/#AM}{Publisher reference}.
\bibitem[BBZ13]{BBZ} Bourgain, Jean; Burq, Nicolas; Zworski, Maciej Control for Schrödinger operators on 2-tori: rough potentials, J. Eur. Math. Soc., Volume 15 (2013) no. 5, pp. 1597-1628 \href{https://ahl.centre-mersenne.org/item/AHL_2019__2__331_0/#BBZ}{Publisher reference}.
\bibitem[BD18]{BD} Bourgain, Jean; Dyatlov, Semyon Spectral gaps without the pressure condition, Ann. Math., Volume 187 (2018) no. 3, pp. 825-867 \href{https://ahl.centre-mersenne.org/item/AHL_2019__2__331_0/#BD}{Publisher reference}.
\bibitem[BG96]{BG} Burq, Nicolas; Gérard, Patrick Condition nécessaire et suffisante pour la contrôlabilité exacte des ondes, C. R. Math. Acad. Sci. Paris, Volume 325 (1996) no. 7, pp. 749-752 \href{https://ahl.centre-mersenne.org/item/AHL_2019__2__331_0/#BG}{Publisher reference}.
\bibitem[BG18]{Bu16} Burq, Nicolas; Gérard, Patrick Stabilisation of wave equations on the torus with rough dampings (2018) (\url{https://arxiv.org/abs/1801.00983}) \href{https://ahl.centre-mersenne.org/item/AHL_2019__2__331_0/#Bu16}{Publisher reference}.
\bibitem[BLR92]{BLR} Bardos, Claude; Lebeau, Gilles; Rauch, Jeffrey Sharp sufficient conditions for the observation, control, and stabilization of waves from the boundary, SIAM J. Control Optimization, Volume 30 (1992) no. 5, pp. 1024-1065 \href{https://ahl.centre-mersenne.org/item/AHL_2019__2__331_0/#BLR}{Publisher reference}.
\bibitem[Bur]{Bu17} Burq, Nicolas Wave control and second-microlocalization on geodesics (in preparation) \href{https://ahl.centre-mersenne.org/item/AHL_2019__2__331_0/#Bu17}{Publisher reference}.
\bibitem[BZ12]{BZ4} Burq, Nicolas; Zworski, Maciej Control for Schrödinger operators on tori, Math. Res. Lett., Volume 19 (2012) no. 2, pp. 309-324 \href{https://ahl.centre-mersenne.org/item/AHL_2019__2__331_0/#BZ4}{Publisher reference}.
\bibitem[DJ18]{DJ} Dyatlov, Semyon; Jin, Long Semiclassical measures on hyperbolic surfaces have full support, Acta Math., Volume 220 (2018) no. 2, pp. 297-339 \href{https://ahl.centre-mersenne.org/item/AHL_2019__2__331_0/#DJ}{Publisher reference}.
\bibitem[Har89]{Ha} Haraux, Alain Séries lacunaires et contrôle semi-interne des vibrations d’une plaque rectangulaire, J. Math. Pures Appl., Volume 68 (1989) no. 4, pp. 457-465 \href{https://ahl.centre-mersenne.org/item/AHL_2019__2__331_0/#Ha}{Publisher reference}.
\bibitem[Ing36]{ing} Ingham, Albert E. Some trigonometrical inequalities with applications to the theory of series, Math. Z., Volume 41 (1936), pp. 367-379 \href{https://ahl.centre-mersenne.org/item/AHL_2019__2__331_0/#ing}{Publisher reference}.
\bibitem[Jaf90]{Ja} Jaffard, Stéphane Contrôle interne exact des vibrations d’une plaque rectangulaire, Port. Math., Volume 47 (1990) no. 4, pp. 423-429 \href{https://ahl.centre-mersenne.org/item/AHL_2019__2__331_0/#Ja}{Publisher reference}.
\bibitem[Jak97]{Jak} Jakobson, Dmitry Quantum limits on flat tori, Ann. Math., Volume 145 (1997) no. 2, pp. 235-266 \href{https://ahl.centre-mersenne.org/item/AHL_2019__2__331_0/#Jak}{Publisher reference}.
\bibitem[Jin17]{Ji} Jin, Long Control for Schrödinger equation on hyperbolic surfaces (2017) (\url{https://arxiv.org/abs/1707.04990,} to appear in Math. Res. Lett.) \href{https://ahl.centre-mersenne.org/item/AHL_2019__2__331_0/#Ji}{Publisher reference}.
\bibitem[Kom92]{Ko} Komornik, Vilmos On the exact internal controllability of a Petrowsky system, J. Math. Pures Appl., Volume 71 (1992) no. 4, pp. 331-342 \href{https://ahl.centre-mersenne.org/item/AHL_2019__2__331_0/#Ko}{Publisher reference}.
\bibitem[Lio88]{Li} Lions, Jacques-Louis Contrôlabilité exacte, perturbations et stabilisation de systèmes distribués. Tome 2: Perturbations, Recherches en Mathématiques Appliquées, 9, Masson, 1988 \href{https://ahl.centre-mersenne.org/item/AHL_2019__2__331_0/#Li}{Publisher reference}.
\bibitem[Mac09]{Ma} Macià, Fabricio Semiclassical measures and the Schrödinger flow on Riemannian manifolds, Nonlinearity, Volume 22 (2009) no. 5, pp. 1003-1020 \href{https://ahl.centre-mersenne.org/item/AHL_2019__2__331_0/#Ma}{Publisher reference}.
\bibitem[Zwo12]{EZB} Zworski, Maciej Semiclassical analysis, Graduate Studies in Mathematics, 138, American Mathematical Society, 2012 \href{https://ahl.centre-mersenne.org/item/AHL_2019__2__331_0/#EZB}{Publisher reference}.
\end{thebibliography}
\end{document}
